\documentclass[11pt]{article}
\usepackage{fontspec,amsmath,amssymb,amsthm,hyperref}
\usepackage[margin=1in]{geometry}
\setlength{\emergencystretch}{3em}
\title{Grönwall comparison with an extra continuous noise input}
\author{LangevinLean mathematical exposition}
\date{}
\begin{document}\maketitle
\noindent Source module: \texttt{ProximalLangevin/MeanFieldForcing.lean}. Compile this proof document with XeLaTeX or LuaLaTeX. The accompanying HTML pairs the mathematical argument with the exact Lean statements and records the scope of the formalization.


In a Banach space let \(F\) be \(L\)-Lipschitz. Compare actual integral equations with a common input \(B\), one extra continuous input \(A\), and a reference forcing \(G\) satisfying a finite drift defect.


\section{Subtract the extra input before applying integral Grönwall}


Suppose \(\|A_t\|\le a\), \(\|F(Y_t)-G_t\|\le C\), and

\[X_t-x=-\int_0^tF(X_s)ds+B_t+A_t,\qquad Y_t-x=-\int_0^tG_sds+B_t.\]

Set \(D_t=X_t-Y_t-A_t\). The common input cancels and \(\|D_t\|\le\int_0^t[L\|D_s\|+La+C]ds\). Grönwall yields \(\|X_t-Y_t\|\le(C+La)\Psi(L,T)+a\) on \([0,T]\). The proof uses continuous integral equations and never differentiates the additional input.




The inputs \(B,A\) may be continuous Brownian paths in \(L^2\); differentiating them would be invalid. Instead set \(D_t=X_t-Y_t-A_t\). Subtracting the two actual integral equations cancels \(B\) and \(A\), leaving \[D_t=-\int_0^t[F(X_s)-G_s]ds,\quad D_0=0.\] Continuity of the integrands proves that this remainder has derivative \(-F(X_t)+G_t\).



If \(F\) is \(L\)-Lipschitz, \(\|A_t\|\le a\) and \(\|F(Y_t)-G_t\|\le C\), then \[\|D'_t\|\le L\|D_t\|+C+La.\] The right-derivative Grönwall theorem gives \(\|D_t\|\le\operatorname{gronwallBound}(0,L,C+La,T)\) uniformly on \([0,T]\). Adding back \(A_t\) produces the final \(+a\). In Lean, \texttt{hDI} identifies the differentiable integral remainder, \texttt{hderiv} proves its derivative, and \texttt{hb} establishes this exact scalar forcing.




\textbf{Correspondence with the displayed declarations.}


\begin{itemize}

\item \texttt{integral\_forcing\_perturbation}: Actual integral equations with common input, an additional continuous input, and a uniformly controlled reference-force defect satisfy this finite-horizon bound. Neither a derivative of the input nor an error recurrence is assumed.
\end{itemize}


\noindent Lean declarations: \texttt{integral\_forcing\_perturbation}.

\section{Factor a common inverse sampling scale}


The explicit zero-initial-error Grönwall function is linear in its forcing coefficient, including its \(L=0\) branch. Thus \(\operatorname{GB}(0,L,C/r,T)=\operatorname{GB}(0,L,C,T)/r\). If \(r>0\), \(\|A_t\|\le a/r\) and \(\|F(Y_t)-G_t\|\le C/r\), substitute into the preceding theorem and factor:

\[\|X_t-Y_t\|\le\frac{(C+La)\Psi(L,T)+a}{r}.\]


The zero-initial-error Grönwall response is linear in its forcing: \(G(0,L,C/r,T)=G(0,L,C,T)/r\), also in the \(L=0\) branch. Apply the preceding theorem with forcing \(C/r\) and extra-input bound \(a/r\), where \(r>0\). Factor \((C+La)/r\) and use that homogeneity to obtain the common \(1/r\) factor in the conclusion. In the particle application \(r=\sqrt N\), so the entire coefficient is independent of \(N\).




\textbf{Correspondence with the displayed declarations.}


\begin{itemize}

\item \texttt{gronwallBound\_zero\_div}: The zero-initial-error Grönwall response is homogeneous in its forcing.
\item \texttt{integral\_forcing\_perturbation\_scaled}: The uniform forcing and input bounds of order \texttt{1/r} give the same scale for the actual solution error, with an explicit coefficient independent of \texttt{r}.
\end{itemize}


\noindent Lean declarations: \texttt{gronwallBound\_zero\_div}, \texttt{integral\_forcing\_perturbation\_scaled}.

\end{document}
